\chapter{Self-Supervised Audio Representations: MERT \emph{(outline)}}\label{ch:selfsup}
\chapterlede{How a network learns what music sounds like with no labels at
all: mask part of the signal, force the model to predict it, and the
representations needed to succeed turn out to be the ones we wanted.}

\begin{incode}
Status: outline. To be written against \texttt{embeddings.py}
(\texttt{m-a-p/MERT-v1-95M}, frame embeddings at $\sim$75 fps, 13 hidden
layers averaged, mean-pooled per moment).
\end{incode}

Planned contents:
\begin{itemize}
  \item \textbf{The masked-prediction template.} BERT-style objectives for
  continuous signals; why predicting raw waveforms fails and targets must be
  discretized or abstracted.
  \item \textbf{Discrete targets via vector quantization.} Codebooks,
  the RVQ (residual VQ) construction used by neural codecs; commitment
  losses; MERT's acoustic teacher.
  \item \textbf{MERT's musical teacher.} The constant-Q transform (CQT):
  logarithmically-spaced filters aligned with equal temperament
  (Appendix~\ref{ch:temperament}); why adding a CQT reconstruction loss
  injects pitch awareness that codec tokens miss.
  \item \textbf{The convolutional front end.} Strided convolutions as
  learned filterbanks; receptive field and frame-rate arithmetic.
  \item \textbf{Which layer knows what.} Probing studies; why the pipeline
  averages all 13 hidden states rather than taking the last; mean-pooling as
  a bag-of-frames assumption and what temporal information it forfeits.
  \item \textbf{From frames to moments.} The pooling map
  $e = \frac{1}{T}\sum_t h_t$ and its effect on the geometry of
  Chapter~\ref{ch:timbre}.
\end{itemize}
